% Part: methods
% Chapter: induction
% Section: inductive-definitions

\documentclass[../../../include/open-logic-section]{subfiles}

\begin{document}

\olfileid{mth}{ind}{idf}

\olsection{تعریف‌های استقرایی}

در منطق، انواع اشیا را بسیار اوقات \emph{به‌طور استقرایی} تعریف
می‌کنیم؛ یعنی قواعدی مشخص می‌کنیم که تعیین می‌کنند چه چیزی شیئی از نوع
مورد تعریف به شمار می‌آید و توضیح می‌دهند چگونه از اشیای پیشینِ آن نوع،
اشیای تازه‌ای از همان نوع بسازیم. برای نمونه، اغلب انواع خاصی از
رشته‌های نمادها، مانند ترم‌ها و !!{formula}های یک زبان، را به‌طور
استقرایی تعریف می‌کنیم. برای مثالی ساده، رشته‌هایی را در نظر بگیرید که
از حروف $\mathrm{a}$، $\mathrm{b}$، $\mathrm{c}$ و $\mathrm{d}$، نماد
~$\circ$ و کروشه‌های $[$ و $]$ تشکیل می‌شوند؛ مانند
«$[[\mathrm{c} \circ \mathrm{d}][$»، «$[\mathrm{a}[]\circ]$»،
«$\mathrm{a}$» یا
«$[[\mathrm{a} \circ \mathrm{b}]\circ \mathrm{d}]$». احتمالاً احساس
می‌کنید در دو رشتهٔ نخست چیزی «نادرست» است: کروشه‌ها در رشتهٔ نخست
اصلاً «موازنه» ندارند، و شاید احساس کنید «$\circ$» باید عبارت‌هایی را
به هم «وصل کند» که خودشان معنادار باشند. رشته‌های سوم و چهارم بهتر به
نظر می‌رسند: به ازای هر «$[$» یک «$]$» پایانی وجود دارد، اگر اصلاً
کروشه‌ای وجود داشته باشد، و برای هر $\circ$ می‌توانیم در دو سوی آن
عبارت‌های «خوش‌ساختی» بیابیم که درون یک جفت کروشه قرار گرفته‌اند.

می‌خواهیم دقیقاً مشخص کنیم چه چیزی «ترم خوش‌ساخت» به شمار می‌آید. نخست،
هر حرف به‌تنهایی خوش‌ساخت است. هر چیزی که فقط یک حرف نباشد باید به شکل
«$[t \circ s]$» باشد، که در آن $s$ و $t$ خود ترم‌هایی خوش‌ساخت‌اند.
برعکس، اگر $t$ و $s$ خوش‌ساخت باشند، می‌توانیم با قراردادن $\circ$
میان آن‌ها و احاطه‌کردنشان با یک جفت کروشه، ترم خوش‌ساخت تازه‌ای بسازیم.
می‌توانیم با این عملیات مجموعهٔ ترم‌های خوش‌ساخت را \emph{تعریف} کنیم.
این یک \emph{تعریف استقرایی} است.

\begin{defn}[ترم‌های خوش‌ساخت]
  مجموعهٔ \emph{ترم‌های خوش‌ساخت} به‌طور استقرایی چنین تعریف می‌شود:
  \begin{enumerate}
  \item هر یک از حروف $\mathrm{a}$، $\mathrm{b}$، $\mathrm{c}$ و
    $\mathrm{d}$ یک ترم خوش‌ساخت است.
  \item اگر $s_1$ و $s_2$ ترم‌هایی خوش‌ساخت باشند، آنگاه
    $[s_1 \circ s_2]$ نیز خوش‌ساخت است.
  \item هیچ چیز دیگری ترم خوش‌ساخت نیست.
  \end{enumerate}
\end{defn}

این تعریف می‌گوید چیزی ترم خوش‌ساخت به شمار می‌آید اگر و تنها اگر
بتوان آن را در شمار متناهی‌ای از گام‌ها بر اساس دو شرط (۱) و (۲) ساخت.
در گام نخست، همهٔ ترم‌های خوش‌ساختی را می‌سازیم که فقط از تک‌حرف‌ها
تشکیل شده‌اند؛ یعنی
\[
\mathrm{a}, \mathrm{b}, \mathrm{c}, \mathrm{d}
\]
در گام دوم، (۲) را بر ترم‌هایی که ساخته‌ایم اعمال می‌کنیم. برای همهٔ
ترکیب‌های دو حرف به این ترم‌ها می‌رسیم:
\[
[\mathrm{a} \circ \mathrm{a}], [\mathrm{a} \circ \mathrm{b}],
[\mathrm{b} \circ \mathrm{a}], \dots, [\mathrm{d} \circ \mathrm{d}]
\]
در گام سوم، (۲) را دوباره بر هر دو ترم خوش‌ساختی که تا آن زمان
ساخته‌ایم اعمال می‌کنیم. ترم‌های خوش‌ساخت تازه‌ای به دست می‌آوریم؛
مانند $[\mathrm{a} \circ [\mathrm{a} \circ \mathrm{a}]]$ که در آن
$t$ همان $\mathrm{a}$ از گام~۱ و $s$ همان
$[\mathrm{a} \circ \mathrm{a}]$ از گام~۲ است، و نیز
$[[\mathrm{b} \circ \mathrm{c}] \circ [\mathrm{d} \circ \mathrm{b}]]$
که از دو ترم $[\mathrm{b} \circ \mathrm{c}]$ و
$[\mathrm{d} \circ \mathrm{b}]$ از گام~۲ ساخته شده است؛ و همین‌طور
ادامه می‌دهیم. بند (۳) مانع آن می‌شود که چیزی که به این روش ساخته نشده
است پنهانی وارد مجموعهٔ ترم‌های خوش‌ساخت شود.

توجه کنید هنوز اثبات نکرده‌ایم که هر رشتهٔ نمادهایی که خوش‌ساخت
«به نظر می‌رسد»، بنا بر این تعریف خوش‌ساخت است. بااین‌حال باید روشن
باشد هر چیزی که می‌توانیم بسازیم واقعاً خوش‌ساخت «به نظر می‌رسد»:
کروشه‌ها موازنه دارند و $\circ$ بخش‌هایی را به هم وصل می‌کند که خود
خوش‌ساخت‌اند.

ویژگی کلیدی تعریف‌های استقرایی این است که اگر بخواهید چیزی دربارهٔ
همهٔ ترم‌های خوش‌ساخت اثبات کنید، تعریف به شما می‌گوید کدام حالت‌ها را
باید بررسی کنید. برای نمونه، اگر گفته شود $t$ ترمی خوش‌ساخت است، تعریف
استقرایی می‌گوید $t$ چه شکلی می‌تواند داشته باشد: $t$ می‌تواند یک حرف
باشد یا برای جفتی از ترم‌های خوش‌ساخت $s_1$ و~$s_2$ به شکل
$[s_1 \circ s_2]$ باشد. به سبب بند (۳)، تنها همین حالت‌ها ممکن‌اند.

هنگام اثبات گزاره‌هایی دربارهٔ همهٔ اعضای مجموعه‌ای که به‌طور استقرایی
تعریف شده است، شکل قوی استقرا اهمیت ویژه‌ای می‌یابد. برای نمونه، فرض
کنید می‌خواهیم اثبات کنیم در هر ترم خوش‌ساخت با طول~$n$، تعداد $[$ها
کمتر از~$n/2$ است. این را می‌توان گزاره‌ای دربارهٔ همهٔ~$n$ها دید:
برای هر $n$، تعداد $[$ها در هر ترم خوش‌ساخت با طول~$n$ کمتر از~$n/2$
است.

\begin{prop}
  برای هر $n$، تعداد $[$ها در یک ترم خوش‌ساخت با طول~$n$ کمتر از
  ~$n/2$ است.
\end{prop}

\begin{proof}
برای اثبات این نتیجه با استقرای قوی، باید صدق گزارهٔ شرطی زیر را نشان
دهیم:
\begin{quote}
  اگر برای هر $l < k$، هر ترم خوش‌ساخت با طول~$l$ کمتر از $l/2$ نماد
  ~$[$ داشته باشد، آنگاه هر ترم خوش‌ساخت با طول~$k$ کمتر از $k/2$ نماد
  ~$[$ دارد.
\end{quote}
برای اثبات این گزارهٔ شرطی، مقدم آن را فرض می‌کنیم؛ یعنی فرض می‌کنیم
برای هر $l<k$، ترم‌های خوش‌ساخت با طول~$l$ کمتر از $l/2$ نماد $[$
دارند. این فرض را فرض استقرا می‌نامیم. می‌خواهیم نشان دهیم همین گزاره
برای ترم‌های خوش‌ساخت با طول~$k$ نیز صادق است.

پس فرض کنید $t$ ترمی خوش‌ساخت با طول~$k$ باشد. چون ترم‌های خوش‌ساخت
به‌طور استقرایی تعریف شده‌اند، دو حالت داریم: (۱)~$t$ خود یک حرف است،
یا (۲)~$t$ برای ترم‌های خوش‌ساخت $s_1$ و~$s_2$ به شکل
$[s_1 \circ s_2]$ است.
\begin{enumerate}
\item $t$ یک حرف است. آنگاه $k = 1$ و تعداد $[$ها در $t$ برابر~$0$
است. چون $0 < 1/2$، گزاره برقرار است.
\item $t$ برای ترم‌های خوش‌ساخت $s_1$ و~$s_2$ به شکل
  $[s_1 \circ s_2]$ است. بگذارید $l_1$ طول~$s_1$ و $l_2$ طول~$s_2$
  باشد. آنگاه طول~$k$ از $t$ برابر $l_1+l_2+3$ است: طول‌های $s_1$ و
  $s_2$ به‌اضافهٔ سه نماد $[$، $\circ$ و $]$. چون $l_1+l_2+3$ همواره
  از $l_1$ بزرگ‌تر است، $l_1 < k$. به همین ترتیب، $l_2 < k$. پس فرض
  استقرا بر ترم‌های $s_1$ و~$s_2$ اعمال می‌شود: تعداد~$m_1$ از $[$ها
  در $s_1$ کمتر از $l_1/2$ است و تعداد~$m_2$ از $[$ها در~$s_2$ کمتر
  از $l_2/2$ است.

  تعداد $[$ها در $t$ برابر است با تعداد $[$ها در~$s_1$، به‌اضافهٔ
  تعداد $[$ها در~$s_2$، به‌اضافهٔ~$1$؛ یعنی $m_1 + m_2 + 1$. چون
  $m_1 < l_1/2$ و $m_2 < l_2/2$، داریم:
  \[
  m_1 + m_2 + 1 < \frac{l_1}{2} + \frac{l_2}{2} + 1 = \frac{l_1+l_2+2}{2} < \frac{l_1+l_2+3}{2} = k/2.
  \]
\end{enumerate}
در هر حالت نشان داده‌ایم تعداد $[$ها در $t$، بر پایهٔ فرض استقرا،
کمتر از $k/2$ است. حکم با استقرای قوی نتیجه می‌شود.
\end{proof}

\begin{prob}
  مجموعهٔ ترم‌های فراخوش‌ساخت را چنین تعریف کنید:
  \begin{enumerate}
  \item هر یک از حروف $\mathrm{a}$، $\mathrm{b}$، $\mathrm{c}$ و
    $\mathrm{d}$ یک ترم فراخوش‌ساخت است.
  \item اگر $s$ ترمی فراخوش‌ساخت باشد، آنگاه $[s]$ نیز فراخوش‌ساخت است.
  \item اگر $s_1$ و $s_2$ ترم‌هایی فراخوش‌ساخت باشند، آنگاه
    $[s_1 \circ s_2]$ نیز فراخوش‌ساخت است.
  \item هیچ چیز دیگری ترم فراخوش‌ساخت نیست.
  \end{enumerate}
  نشان دهید تعداد $[$ها در یک ترم فراخوش‌ساخت~$t$ با طول~$n$ حداکثر
  $n/2 +1$ است.
\end{prob}

\end{document}
